\scanpage{63}
\mspage{106}
\begin{flushright}number 371\end{flushright}

$F=\handwritten{\prod}_{n}E_n$. Let us show that one can define on $F$ a ring structure which induces on $E$ the given ring structure ; indeed, every element of $F$ is of the form $(u_n)$, where $u_n\handwritten{\in}E_n$ ; one defines on $F$ a multiplication by setting $(u_n)(v_n)=(w_n)$, with
\[
w_n=\handwritten{\sum}_{p+q=n}u_pv_q
\tag*{(28 bis)}
\]
which makes sense for every $n$, since the right-hand side has only finitely many terms ; the element $1\handwritten{\in}E_{\unsure{o|0}}$ of $E$ is clearly an identity element for this law and associativity is checked immediately ; one thus has on $E$ a ring structure (commutative if $E$ is commutative). One can likewise define on $B=\handwritten{\prod}_{n}A_n$ a commutative ring structure ; finally, one defines on $F$ a $B$-\uline{algebra} structure in an analogous way : if $(\handwritten{\alpha}_n)$ is an element of $B$, one will set $(\handwritten{\alpha}_n)(u_n)=(w'_n)$ with
\[
w'_n=\handwritten{\sum}_{p+q=n}\handwritten{\alpha}_pu_q.
\tag*{(28 ter)}
\]

We shall say that $B$ (resp. $F$) is the \uline{canonical completion} of the graded ring $A$ (resp. of the graded $A$-algebra $E$).

\begin{quote}
\handwritten{*}This terminology is justified by considering on $B$ (resp. $E$) the \uline{product} topology of the discrete topologies on the $A_n$ (resp. $E_n$) ; then one sees at once that multiplication in $F$ (resp. the external $B$-module law on $F$) is the \uline{extension by continuity} to $F\handwritten{\times}F$ (resp. $B\handwritten{\times}F$) of the multiplication $E\handwritten{\times}E\handwritten{\to}E$ (resp. of the external law $A\handwritten{\times}E\handwritten{\to}E$).\handwritten{*}
\end{quote}

Let $E'$ be a second graded $A$-algebra of type $Z$ with positive degrees, and let $F'$ be its canonical completion. Every homomorphism $f:E\handwritten{\to}E'$ of graded $A$-algebras then extends canonically to a homomorphism $g:F\handwritten{\to}F'$ of $B$-algebras ; it suffices to set, for every $u=(u_n)\handwritten{\in}F$
\[
g(u)=(f(u_n))
\]
and one checks at once that this does indeed define a homomorphism of $B$-algebras.
% archive sheet number: 62
